\begin{definitionbox}{Definition}
\begin{definition}[Point-slope form]
\label{def:point-slope-form}
Let $P_1=(x_1,y_1)$ be a point in the Euclidean plane and let
$m\in\mathbb{R}$. The point-slope form of the line through $P_1$ with slope
$m$ is
\[
  y-y_1=m(x-x_1).
\]
\end{definition}
\end{definitionbox}

\begin{dependencies}
\begin{itemize}
  \item \hyperref[def:reals]{Real Numbers}
  \item \hyperref[def:cartesian-product]{Cartesian Product}
  \item \hyperref[def:subtraction-on-r]{Subtraction on $\mathbb{R}$}
  \item \hyperref[def:multiplication-on-r]{Multiplication on $\mathbb{R}$}
\end{itemize}
\end{dependencies}

\begin{remark*}[Standard quantified statement]
For every point $P_1=(x_1,y_1)\in\mathbb{R}^2$ and every
$m\in\mathbb{R}$, the corresponding line is the set of all
$(x,y)\in\mathbb{R}^2$ satisfying $y-y_1=m(x-x_1)$.
\end{remark*}

\begin{remark*}[Derivation]
The formula records the condition that the slope from $P_1$ to an arbitrary
point $(x,y)$ on the line is constant and equal to $m$.
\end{remark*}

\begin{definitionbox}{Definition}
\begin{definition}[Two-point form]
\label{def:two-point-form}
Let $P_1=(x_1,y_1)$ and $P_2=(x_2,y_2)$ be distinct points with
$x_1\neq x_2$. The two-point form of the line through $P_1$ and $P_2$ is
\[
  y-y_1=\left(\frac{y_2-y_1}{x_2-x_1}\right)(x-x_1).
\]
\end{definition}
\end{definitionbox}

\begin{dependencies}
\begin{itemize}
  \item \hyperref[def:point-slope-form]{Point-Slope Form}
  \item \hyperref[def:subtraction-on-r]{Subtraction on $\mathbb{R}$}
  \item \hyperref[thm:r-is-a-field]{$\mathbb{R}$ Is a Field}
\end{itemize}
\end{dependencies}

\begin{remark*}[Standard quantified statement]
For all distinct points $P_1=(x_1,y_1)$ and $P_2=(x_2,y_2)$ in
$\mathbb{R}^2$ with $x_1\neq x_2$, the line through $P_1$ and $P_2$ is the set
of all $(x,y)\in\mathbb{R}^2$ satisfying the displayed two-point equation.
\end{remark*}

\begin{definitionbox}{Definition}
\begin{definition}[General linear form]
\label{def:general-linear-form}
The general linear form of a line in $\mathbb{R}^2$ is
\[
  Ax+By+C=0,
\]
where $A,B,C\in\mathbb{R}$ and $A^2+B^2\neq 0$.
\end{definition}
\end{definitionbox}

\begin{dependencies}
\begin{itemize}
  \item \hyperref[def:reals]{Real Numbers}
  \item \hyperref[def:addition-on-r]{Addition on $\mathbb{R}$}
  \item \hyperref[def:multiplication-on-r]{Multiplication on $\mathbb{R}$}
  \item \hyperref[thm:r-is-a-field]{$\mathbb{R}$ Is a Field}
\end{itemize}
\end{dependencies}

\begin{remark*}[Standard quantified statement]
For all $A,B,C\in\mathbb{R}$ with $A^2+B^2\neq 0$, the equation
$Ax+By+C=0$ determines a nondegenerate line in $\mathbb{R}^2$.
\end{remark*}

\begin{remark*}[Interpretation]
The vector $(A,B)$ is normal to the line. This form includes vertical lines
and therefore avoids the exceptional cases present in slope-based forms.
\end{remark*}

\begin{definitionbox}{Definition}
\begin{definition}[Vector parametric form]
\label{def:vector-parametric-form}
A line in $\mathbb{R}^2$ may be represented as the image of
$\mathbf r:\mathbb{R}\to\mathbb{R}^2$ given by
\[
  \mathbf r(t)=\mathbf r_0+t\mathbf v,
\]
where $\mathbf r_0\in\mathbb{R}^2$ is a point on the line and
$\mathbf v\in\mathbb{R}^2$ is a nonzero direction vector.
\end{definition}
\end{definitionbox}

\begin{dependencies}
\begin{itemize}
  \item \hyperref[def:reals]{Real Numbers}
  \item \hyperref[def:cartesian-product]{Cartesian Product}
  \item \hyperref[def:function]{Function}
  \item \hyperref[def:addition-on-r]{Addition on $\mathbb{R}$}
  \item \hyperref[def:multiplication-on-r]{Multiplication on $\mathbb{R}$}
\end{itemize}
\end{dependencies}

\begin{remark*}[Standard quantified statement]
For every $\mathbf r_0\in\mathbb{R}^2$ and every nonzero
$\mathbf v\in\mathbb{R}^2$, the set
$\{\mathbf r_0+t\mathbf v:t\in\mathbb{R}\}$ is a line in $\mathbb{R}^2$.
\end{remark*}

\begin{definitionbox}{Definition}
\begin{definition}[Intercept form]
\label{def:intercept-form}
If a line meets the $x$-axis at $(a,0)$ and the $y$-axis at $(0,b)$ with
$a,b\neq 0$, its intercept form is
\[
  \frac{x}{a}+\frac{y}{b}=1.
\]
\end{definition}
\end{definitionbox}

\begin{dependencies}
\begin{itemize}
  \item \hyperref[def:two-point-form]{Two-Point Form}
  \item \hyperref[thm:r-is-a-field]{$\mathbb{R}$ Is a Field}
\end{itemize}
\end{dependencies}

\begin{remark*}[Standard quantified statement]
For all nonzero $a,b\in\mathbb{R}$, the line with intercepts $(a,0)$ and
$(0,b)$ is the set of all $(x,y)\in\mathbb{R}^2$ satisfying
$x/a+y/b=1$.
\end{remark*}

\begin{remark*}[Derivation]
The intercept form follows by applying the two-point form to the points
$(a,0)$ and $(0,b)$ and then rearranging the resulting equation.
\end{remark*}

\begin{definitionbox}{Definition}
\begin{definition}[Affine parametric line in $\mathbb{R}^n$]
\label{def:affine-parametric-line}
Given distinct points $\mathbf x_1,\mathbf x_2\in\mathbb{R}^n$, the line
through them is
\[
  \mathbf x(t)=\mathbf x_1+t(\mathbf x_2-\mathbf x_1),
  \qquad t\in\mathbb{R}.
\]
\end{definition}
\end{definitionbox}

\begin{dependencies}
\begin{itemize}
  \item \hyperref[def:reals]{Real Numbers}
  \item \hyperref[def:cartesian-product]{Cartesian Product}
  \item \hyperref[def:subtraction-on-r]{Subtraction on $\mathbb{R}$}
  \item \hyperref[def:addition-on-r]{Addition on $\mathbb{R}$}
  \item \hyperref[def:multiplication-on-r]{Multiplication on $\mathbb{R}$}
\end{itemize}
\end{dependencies}

\begin{remark*}[Standard quantified statement]
For all distinct $\mathbf x_1,\mathbf x_2\in\mathbb{R}^n$, the set
$\{\mathbf x_1+t(\mathbf x_2-\mathbf x_1):t\in\mathbb{R}\}$ is the affine
line through $\mathbf x_1$ and $\mathbf x_2$.
\end{remark*}

\begin{remark*}[Interpretation]
For $0\le t\le 1$, the same formula describes the line segment joining
$\mathbf x_1$ and $\mathbf x_2$. For arbitrary real $t$, it describes the
entire affine line.
\end{remark*}

\begin{definitionbox}{Definition}
\begin{definition}[Normal form]
\label{def:normal-form-line}
A line in $\mathbb{R}^2$ may be written in normal form as
\[
  x\cos\alpha+y\sin\alpha-p=0,
\]
where $p\ge 0$ is the distance from the origin to the line.
\end{definition}
\end{definitionbox}

\begin{dependencies}
\begin{itemize}
  \item \hyperref[def:general-linear-form]{General Linear Form}
  \item \hyperref[def:sine-cosine]{Sine and Cosine}
  \item \hyperref[def:order-on-r]{Order on $\mathbb{R}$}
\end{itemize}
\end{dependencies}

\begin{remark*}[Standard quantified statement]
For every $p\ge 0$ and every real angle $\alpha$, the equation
$x\cos\alpha+y\sin\alpha-p=0$ describes a line whose unit normal direction is
$(\cos\alpha,\sin\alpha)$.
\end{remark*}

\begin{definitionbox}{Definition}
\begin{definition}[Symmetric form]
\label{def:symmetric-form-line}
Given a point $(x_1,y_1)$ and direction cosines
$(\cos\theta,\sin\theta)$, the symmetric form of a line is
\[
  \frac{x-x_1}{\cos\theta}
  =
  \frac{y-y_1}{\sin\theta}
  =
  r.
\]
\end{definition}
\end{definitionbox}

\begin{dependencies}
\begin{itemize}
  \item \hyperref[def:vector-parametric-form]{Vector Parametric Form}
  \item \hyperref[def:sine-cosine]{Sine and Cosine}
  \item \hyperref[thm:r-is-a-field]{$\mathbb{R}$ Is a Field}
\end{itemize}
\end{dependencies}

\begin{remark*}[Standard quantified statement]
For every point $(x_1,y_1)$ and every direction with nonzero displayed
denominators, the symmetric equation describes the same line as the
parametric equations
$x=x_1+r\cos\theta$ and $y=y_1+r\sin\theta$.
\end{remark*}
